\section{实验}
\label{sec:experiments}

\subsection{实验设置}

\paragraph{数据集。}
我们使用当前 MiniOneRec 复现主线中已经准备好的 Industrial 和 Office Amazon 数据集。Industrial 包含 $36{,}259$ 条训练交互、$4{,}533$ 条测试交互和 $3{,}686$ 个 item；Office 包含 $38{,}924$ 条训练交互、$4{,}866$ 条测试交互和 $3{,}459$ 个 item。两者都沿用仓库中的时间划分。

\paragraph{Tokenizer backbone。}
除非特别说明，本文所有训练时 tokenizer 比较都使用与上游 MiniOneRec 对齐的 \rqvae{} 训练制度：三层量化器、每层 codebook 大小为 $256$、量化 embedding 维度为 $32$、AdamW、学习率 $10^{-3}$、训练 $10{,}000$ epochs、batch size 为 $20{,}480$。

\paragraph{评估指标。}
我们报告三类指标：
\begin{itemize}[leftmargin=1.4em]
  \item \textbf{Probe 指标}：不同 ambiguity bucket 上的 target-better rate，重点关注 same-$l_2$；
  \item \textbf{Tokenizer 指标}：最佳 train-stage collision、最终 \texttt{sid-generate} collision，以及最大冲突大小；
  \item \textbf{局部歧义指标}：平均目标 $l_2$ 叶子数、目标落入多叶 same-$l_2$ bucket 的比例、目标落入深拥挤 same-$l_2$ bucket 的比例，以及在固定前缀下第三级 code 的条件熵。
\end{itemize}

\subsection{设计空间 probe：spectral middle view 的必要性}

我们首先做不改 tokenizer 本体的 graph probe。给定一个现有 top-1 SID 错误样本，我们比较某个 graph view 是否更偏向真实 target，而不是当前被错误偏好的候选。这样可以在接入训练之前，直接检查不同图视图对 local ambiguity 的判别力。

\begin{table}[t]
  \centering
  \caption{same-$l_2$ bucket 上的最佳 middle-view probe 结果。spectral middle view 明显强于第一版 heuristic middle view，这也是 MGR-SID v1 采用 spectral $G_{\mathrm{mid}}$ 的直接原因。}
  \label{tab:probe}
  \begin{tabular}{lccc}
    \toprule
    数据集 & 最佳 heuristic middle view & 最佳 spectral middle view & spectral 视图名称 \\
    \midrule
    Industrial & 0.2747 & 0.5401 & \texttt{fagsp\_mid\_base} \\
    Office & 0.2381 & 0.7619 & \texttt{fagsp\_mid\_base} \\
    \bottomrule
  \end{tabular}
\end{table}

\Cref{tab:probe} 说明，真正值得被显式设计的是 middle-resolution graph structure 本身。对 same-$l_2$ 歧义而言，最佳 spectral middle view 在 Industrial 上几乎把 heuristic middle view 的判别力提升到了两倍，在 Office 上甚至超过三倍。这一结果直接推动我们把 $G_{\mathrm{mid}}$ 从 heuristic proxy 升级为 spectral graph。

\subsection{在 upstream-aligned tokenizer 训练上的正式比较}

接下来我们测试图信号进入 tokenizer 学习后是否仍然成立。这里所有方法都使用 \Cref{tab:provenance} 中已经验证过的 upstream-aligned 训练制度。我们比较 semantic baseline、uniform graph regularization，以及 hierarchy-aware 的 MGR-SID v1。

\begin{table}[t]
  \centering
  \caption{Industrial 上的 tokenizer 比较结果。只看 \texttt{sid-train} 会得出误导性结论；真正的排序应当看最终 \texttt{sid-generate} 后的 SID。}
  \label{tab:main}
  \begin{tabular}{lccc}
    \toprule
    方法 & 最佳 train collision & 最终 index collision & Max conflict \\
    \midrule
    Semantic baseline & 0.1066 & 0.00353 & 2 \\
    Uniform graph regularization & 0.2458 & 0.00597 & 6 \\
    Hierarchy-aware MGR-SID v1 & 0.1319 & \textbf{0.00326} & 2 \\
    \bottomrule
  \end{tabular}
\end{table}

\Cref{tab:main} 给出两个重要事实。第一，uniform graph regularization 明显有害，说明“把 graph 加进去”本身并不够。第二，train-stage collision 与最终 tokenizer 质量并不一致。虽然 semantic baseline 在 raw \texttt{sid-train} collision 上最好，但经过官方 \texttt{sid-generate} 后，hierarchy-aware tokenizer 反而得到最好的最终 SID。这也再次说明，tokenizer 比较必须落在最终导出的 index 上。

\subsection{最终 tokenizer 是否真的降低了 local ambiguity}

仅有更低的 final collision 还不够，我们还需要确认它是否真的改善了我们最关心的失败模式。为此，我们直接比较 baseline final index 与 hierarchy final index 在 same-$l_2$ 层面的结构差异。

\begin{table}[t]
  \centering
  \caption{Industrial 上最终 SID 的局部歧义比较。Hierarchy-aware tokenizer 在 catalog 级和 test-weighted 两种口径下都降低了 same-$l_2$ 的拥挤程度。}
  \label{tab:ambiguity}
  \begin{tabular}{lccc}
    \toprule
    指标 & Baseline & Hierarchy & Delta \\
    \midrule
    Mean target $l_2$ leaf count & 4.7999 & 4.4498 & -0.3501 \\
    Fraction targets in multi-leaf $l_2$ & 0.6894 & 0.6131 & -0.0763 \\
    Fraction targets in $l_2$ with $\ge 4$ leaves & 0.3283 & 0.2828 & -0.0454 \\
    Mean target level-3 entropy under $l_2$ & 1.4533 & 1.2935 & -0.1598 \\
    \bottomrule
  \end{tabular}
\end{table}

\Cref{tab:ambiguity} 表明，这不是简单的 code re-permutation。Hierarchy-aware tokenizer 的确让 same-$l_2$ bucket 更干净：目标 item 平均面对的 $l_2$ 叶子数更少，落在多叶与深拥挤 same-$l_2$ bucket 中的比例更低，同时给定 prefix 后第三级 code 的条件熵也更低。在更细的 movement analysis 中，有 $17.32\%$ 的测试目标从多叶 same-$l_2$ bucket 迁移到单叶 bucket，而只有 $9.68\%$ 迁移到更拥挤的方向。这与我们最初关于 local leaf ambiguity 的诊断完全一致。

\subsection{第一轮下游验证：v1 只实现了部分传递}

我们首先把 \Cref{tab:main,tab:ambiguity} 中的两套最终 Industrial tokenizer 推进到下游 SFT。这里使用的是一组 \emph{internal control}：upstream-aligned semantic control 与 hierarchy-aware v1。在这一步里，我们刻意保持整个管线尽量稳定，只替换 tokenizer，以便先观察 tokenizer 改动是否会传到下游行为。

\begin{table}[t]
  \centering
  \caption{Industrial 上 internal control 的 SFT 结果。MGR-SID v1 在短程排序上有正信号，但还没有形成广义的 $@10+$ 增益。}
  \label{tab:sft_v1}
  \begin{tabular}{lcccccc}
    \toprule
    运行 & N@1 & N@3 & N@10 & H@1 & H@3 & H@10 \\
    \midrule
    internal semantic control & 0.0631 & 0.0777 & 0.0943 & 0.0631 & 0.0882 & 0.1343 \\
    internal hierarchy-aware v1 & 0.0627 & 0.0802 & 0.0936 & 0.0627 & 0.0929 & 0.1304 \\
    \bottomrule
  \end{tabular}
\end{table}

\Cref{tab:sft_v1} 说明，tokenizer 侧收益确实能传到下游，但只是部分传递：相对于 internal semantic control，v1 提升了 $\mathrm{NDCG}@3$ 与 $\mathrm{HR}@3$，在 $@5$ 上也有轻微增益，但在 $@10+$ 上仍然落后。正是这组结果促使我们把重点从“继续换 graph carrier”转向“重做 tokenizer 训练接口”。

\subsection{Ambiguity-aware tokenizer v2}

v1 的 top-$k$ 与错误分布分析表明：graph regularization 在拥挤局部歧义样本上有帮助，但会过度修正那些本来已经稳定的语义区域。因此我们提出 \texttt{v2}，核心改动是：
\begin{itemize}[leftmargin=1.4em]
  \item 用由 semantic density、semantic-collaborative disagreement 与 graph competition 组成的 \emph{offline ambiguity prior} 去调节 graph supervision；
  \item 加入 \emph{semantic-structure retention}，在低歧义区域保护原始语义几何。
\end{itemize}

第一版 \texttt{v2} 只采用离线 combined proxy，因为当前的 online uncertainty proxy 在 sanity check 中还不够稳定。

\begin{table}[t]
  \centering
  \caption{引入 ambiguity-aware supervision 后的 tokenizer 比较。v2 没有进一步降低 final collision，但得到的是最干净的局部层级结构。}
  \label{tab:v2tok}
  \begin{tabular}{lcccc}
    \toprule
    方法 & Final collision & Mean target $l_2$ leaves & Multi-leaf same-$l_2$ & Mean level-3 entropy \\
    \midrule
    internal semantic control & 0.00353 & 4.7999 & 0.6894 & 1.4533 \\
    hierarchy-aware v1 & \textbf{0.00326} & 4.4498 & 0.6131 & 1.2935 \\
    ambiguity-aware v2 & 0.00353 & \textbf{4.3422} & \textbf{0.4873} & \textbf{1.1001} \\
    \bottomrule
  \end{tabular}
\end{table}

\Cref{tab:v2tok} 的意义在于：collision 并不是 tokenizer 的全部。虽然 \texttt{v2} 在 final collision 上只打平了 internal semantic control，但它得到的是目前最干净的 local hierarchy。这说明 ambiguity-aware 设计并不是“再把 graph 加重”，而是更有选择地把 graph supervision 放在真正高歧义的区域上。

\subsection{在 recipe-aligned 的原版 MiniOneRec 上做下游比较}

接下来我们先把 \texttt{v2} 推进到真正更公平的下游对照中：与 \emph{原版 MiniOneRec} 在相同 task recipe 下比较。第一轮 \texttt{v2} 使用的 task recipe 为：
\begin{itemize}[leftmargin=1.4em]
  \item \texttt{title\_history2sid = on}
  \item \texttt{SID-description alignment = off}
\end{itemize}

因此，最公平的原版 MiniOneRec 对照是使用相同 task recipe 的历史运行。为了给出仓库内部的整体位置，我们同时报告最强的原版 MiniOneRec SFT 与 RL 结果。

\begin{table}[t]
  \centering
  \caption{Industrial 上 \texttt{v2} 的下游比较。v2 已经超过 recipe-aligned 的原版 MiniOneRec，但还没有超过最强的原版 MiniOneRec recipe。}
  \label{tab:v2sft}
  \begin{tabular}{lcccccc}
    \toprule
    运行 & N@1 & N@3 & N@10 & H@1 & H@3 & H@10 \\
    \midrule
    recipe-aligned original MiniOneRec & 0.0624 & 0.0798 & 0.0987 & 0.0624 & 0.0929 & \textbf{0.1471} \\
    strongest original MiniOneRec SFT & 0.0671 & \textbf{0.0850} & \textbf{0.1037} & 0.0671 & \textbf{0.0984} & 0.1509 \\
    strongest original MiniOneRec RL & \textbf{0.0732} & 0.0890 & 0.1073 & \textbf{0.0732} & 0.1004 & 0.1513 \\
    ambiguity-aware v2 & 0.0704 & 0.0839 & 0.1008 & 0.0704 & 0.0942 & 0.1425 \\
    \bottomrule
  \end{tabular}
\end{table}

\Cref{tab:v2sft} 更清楚地说明了当前方法的真实位置。相对于 \emph{recipe-aligned 的原版 MiniOneRec}，\texttt{v2} 在所有报告的 NDCG cutoff 上都更好：
\[
\Delta \mathrm{NDCG}@1 = +0.0079,\;
\Delta \mathrm{NDCG}@3 = +0.0041,\;
\Delta \mathrm{NDCG}@5 = +0.0036,\;
\Delta \mathrm{NDCG}@10 = +0.0021.
\]
同时它也提升了 $\mathrm{HR}@1$、$\mathrm{HR}@3$ 与 $\mathrm{HR}@5$，只是 \(\mathrm{HR}@10\) 仍略低。这说明 \texttt{v2} 已经不是只在 internal control 上成立的方法，它已经在与原版 MiniOneRec 的公平对照中表现出下游增益。

但这张表也提醒我们：\texttt{v2} 还没有超过最强的原版 MiniOneRec recipe。当前最强的原版 MiniOneRec SFT 仍然是 \texttt{title\_history2sid\_off + desc\_align\_p05}，而 strongest RL 仍然更强。到这一步为止，我们知道 \texttt{v2} 在公平 task recipe 下已经成立，但 strongest recipe mismatch 的具体来源还没有拆清。

\subsection{Recipe isolation：真正的冲突来自 \texttt{title\_history2sid\_off}}

为了弄清 strongest original recipe 为什么不能直接迁移到 \texttt{v2} 上，我们补齐了 \texttt{v2} 的四格 recipe matrix：
\begin{itemize}[leftmargin=1.4em]
  \item \texttt{title\_history2sid\_on + desc\_align\_off}
  \item \texttt{title\_history2sid\_on + desc\_align\_p05}
  \item \texttt{title\_history2sid\_off + desc\_align\_off}
  \item \texttt{title\_history2sid\_off + desc\_align\_p05}
\end{itemize}

\begin{table}[t]
  \centering
  \caption{Industrial 上 \texttt{v2} 的 recipe isolation 结果。当前最优下游 recipe 是 \texttt{title\_history2sid\_on + desc\_align\_p05}，主要负因子是关闭 \texttt{title\_history2sid}。}
  \label{tab:v2recipe}
  \begin{tabular}{lcccccc}
    \toprule
    运行 & N@1 & N@3 & N@10 & H@1 & H@3 & H@10 \\
    \midrule
    \texttt{v2\_on\_off} & 0.0704 & 0.0839 & 0.1008 & 0.0704 & 0.0942 & 0.1425 \\
    \texttt{v2\_on\_p05} & \textbf{0.0706} & \textbf{0.0845} & \textbf{0.1027} & \textbf{0.0706} & \textbf{0.0951} & \textbf{0.1463} \\
    \texttt{v2\_off\_off} & 0.0589 & 0.0741 & 0.0913 & 0.0589 & 0.0854 & 0.1339 \\
    \texttt{v2\_off\_p05} & 0.0591 & 0.0737 & 0.0899 & 0.0591 & 0.0843 & 0.1308 \\
    \bottomrule
  \end{tabular}
\end{table}

\Cref{tab:v2recipe} 直接拆清了 strongest-recipe mismatch。只要 \texttt{title\_history2sid} 保持开启，\texttt{desc\_align\_p05} 对 \texttt{v2} 是净正增益：
\[
\Delta \mathrm{NDCG}@10 = +0.0019,\qquad
\Delta \mathrm{HR}@10 = +0.0038.
\]
相反，只要把 \texttt{title\_history2sid} 关掉，无论 \texttt{desc\_align} 是否开启，性能都会显著下滑。在 \texttt{desc\_align\_off} 下：
\[
\Delta \mathrm{NDCG}@10 = -0.0096,\qquad
\Delta \mathrm{HR}@10 = -0.0086,
\]
而在 \texttt{desc\_align\_p05} 下，下滑更大：
\[
\Delta \mathrm{NDCG}@10 = -0.0128,\qquad
\Delta \mathrm{HR}@10 = -0.0154.
\]

因此，强 recipe 失配的主因并不是 \texttt{desc\_align\_p05}，而是 \texttt{title\_history2sid\_off}。对于当前 \texttt{v2} 来说，最佳的下游 recipe 已经更新为 \texttt{title\_history2sid\_on + desc\_align\_p05}。

我们进一步把这条最优 \texttt{v2} 设置与 strongest original MiniOneRec SFT 做了更细的比较。两者的差距现在已经很小：
\[
\Delta \mathrm{NDCG}@10 = -0.0010,\qquad
\Delta \mathrm{HR}@10 = -0.0046.
\]
更细的 top-$k$ 分析说明，\texttt{v2\_on\_p05} 已经在 rank-1 上更强，same-prefix top-1 错误也明显更少；strongest original SFT 的主要优势则集中在 \texttt{top3/top5/top10} 的中段 beam retention。因此，当前剩余问题已经不再像 tokenizer 方向错误，而更像是 downstream interaction gap。

\subsection{当前结果的解释}

综合起来，目前实验支持的是一个更准确的结论：
\begin{enumerate}[leftmargin=1.4em]
  \item 公平的 tokenizer 比较必须建立在与上游 MiniOneRec 对齐的训练制度上；
  \item graph-structured collaborative information 对 local same-$l_2$ disambiguation 确实有信号，而且其中最关键的载体是 spectral middle-resolution graph；
  \item 仅有 uniform hierarchy-aware 图注入还不够，因为它会过度修正已经稳定的语义区域；
  \item ambiguity-aware graph supervision 加 semantic-structure retention 能得到比 internal control 与 v1 都更干净的局部层级结构；
  \item 这种更干净的 tokenizer 已经能在 recipe-aligned 的原版 MiniOneRec 上带来下游增益；
  \item 经过 recipe isolation 之后，当前 \texttt{v2} 的最佳下游 setting 已经明确为 \texttt{title\_history2sid\_on + desc\_align\_p05}；
  \item strongest original MiniOneRec recipe 仍然略强，但剩余差距已经收缩成一个中段 beam retention 问题，而不是 tokenizer 侧的彻底失败。
\end{enumerate}

因此，当前阶段最重要的结论不是“我们已经赢了整个系统”，而是：
\begin{quote}
tokenizer 侧方法现在已经同时得到结构证据、recipe-aligned 下游证据以及 recipe-isolation 证据的支持。当前最佳 \texttt{v2} setting 已经非常接近 strongest original MiniOneRec SFT，剩余问题主要是 downstream retention，而不是 ambiguity-aware tokenizer 本身的方向错误。
\end{quote}
